\chapter{Kemagnetan dan Medan Magnet}\label{ch:g11:magnetic-fields}

Para pelaut mengemudikan kapalnya dengan jarum bermagnet selama seribu
tahun sebelum ada yang tahu mengapa jarum itu menunjuk ke utara.
Jawabannya tiba pada tahun 1820, ketika sebuah arus membuat kompas
tersentak: kemagnetan dibuat oleh muatan yang bergerak. Bab ini
memetakan \omterm{def:g11:magnetic-fields:bfield}{medan magnet} --- milik \omterm{def:g11:magnetic-fields:magnet}{magnet}, milik Bumi, milik arus ---
lalu menutup untainya: \omterm{def:g11:electric-gravitational-fields:field}{medan} itu mendorong balik arusnya, dan begitulah
setiap motor listrik berputar.

\section{Magnet dan medan magnet}

\begin{definition}[Magnet dan kutubnya]\label{def:g11:magnetic-fields:magnet}
Sebuah \emph{magnet}\index{magnet} menarik besi dan, ketika digantung
pada seutas benang, berputar menghadap utara. Pengaruhnya terpusat pada
dua \emph{kutub magnet}\index{kutub magnet}, yaitu kutub \emph{utara}
(ujung yang mencari utara geografis) dan kutub \emph{selatan}: kutub
sejenis tolak-menolak, kutub tak sejenis tarik-menarik. Menggergaji
sebuah magnet menjadi dua menghasilkan dua magnet yang lengkap ---
belum pernah ada percobaan yang mengasingkan satu kutub.
\end{definition}

\begin{definition}[Medan magnet]\label{def:g11:magnetic-fields:bfield}
Sebuah \omterm{def:g11:magnetic-fields:magnet}{magnet} mengubah ruang di sekelilingnya
(\cref{ch:g11:electric-gravitational-fields}): \omterm{def:g11:magnetic-fields:magnet}{magnet} itu menciptakan
\emph{medan magnet}\index{medan magnet} $\vect B$ di mana pun. Kompas
membaca arahnya (jarumnya berbaris searah $\vect B$, dengan ujung
utaranya di depan); \emph{penduga Hall}\index{penduga Hall} mengukur
besarnya, dalam \emph{tesla}\index{tesla} (\unit{T}).
\end{definition}

\begin{definition}[Garis medan magnet]\label{def:g11:magnetic-fields:lines}
\emph{Garis medan magnet}\index{garis medan magnet} adalah kurva yang
menyinggung $\vect B$, dan berdesakan di tempat \omterm{def:g11:electric-gravitational-fields:field}{medannya} kuat. Di luar
sebuah \omterm{def:g11:magnetic-fields:magnet}{magnet}, garis itu berjalan dari kutub utara ke kutub selatan;
tidak seperti garis \omterm{def:g11:electric-gravitational-fields:efield}{medan listrik}, garis itu tidak berujung:
masing-masing menutup melalui tubuh \omterm{def:g11:magnetic-fields:magnet}{magnetnya}.
\end{definition}

\begin{omfigure}
\begin{tikzpicture}[scale=0.9,
  fline/.style={omThm, postaction={decorate},
    decoration={markings, mark=at position 0.5 with {\arrow{Stealth}}}}]
  \fill[omProp!15] (-1.4,-0.35) rectangle (0,0.35);
  \fill[omDef!15] (0,-0.35) rectangle (1.4,0.35);
  \draw[thick] (-1.4,-0.35) rectangle (1.4,0.35);
  \draw (0,-0.35) -- (0,0.35);
  \node[font=\small] at (-0.7,0) {S};
  \node[font=\small] at (0.7,0) {U};
  \draw[fline] (1.4,0.2) .. controls (2.6,0.9) and (-2.6,0.9) .. (-1.4,0.2);
  \draw[fline] (1.4,0.3) .. controls (3.8,2.0) and (-3.8,2.0) .. (-1.4,0.3);
  \draw[fline] (1.4,-0.2) .. controls (2.6,-0.9) and (-2.6,-0.9) .. (-1.4,-0.2);
  \draw[fline] (1.4,-0.3) .. controls (3.8,-2.0) and (-3.8,-2.0) .. (-1.4,-0.3);
  \draw[fline] (1.4,0) -- (2.95,0);
  \draw[omExo, thick] (0,1.1) circle (0.22);
  \draw[-Stealth, omDef, thick] (0.22,1.1) -- (-0.22,1.1);
  \draw[omExo, thick] (2.45,0) circle (0.22);
  \draw[-Stealth, omDef, thick] (2.23,0) -- (2.67,0);
\end{tikzpicture}

{\small \omterm{def:g11:electric-gravitational-fields:lines}{Garis medan} sebuah \omterm{def:g11:magnetic-fields:magnet}{magnet} batang: keluar dari U, masuk ke S,
dan menutup melalui tubuhnya. Kompas yang ditempatkan di mana pun
berhenti sejajar garis yang melewatinya.}
\end{omfigure}

\section{Bumi adalah sebuah magnet}

\begin{definition}[Medan magnet Bumi]\label{def:g11:magnetic-fields:earth}
Arus di dalam inti besi cair Bumi menjadikan planet itu sebuah \omterm{def:g11:magnetic-fields:magnet}{magnet}
raksasa: di permukaannya $B \approx \qty{50}{\micro T}$, dan jarum yang
bebas berbaris searah komponen mendatar \omterm{def:g11:electric-gravitational-fields:field}{medannya} --- itulah kompas.
Kutub geomagnetnya bukan kutub geografisnya; sudut antara utara kompas
dan utara sejati disebut \emph{deklinasi}\index{deklinasi}, yang
dikoreksi pada setiap penunjukan arah.
\end{definition}

\begin{example}[Mengoreksi sebuah arah]\label{ex:g11:magnetic-fields:bearing}
Di tempat \omterm{def:g11:magnetic-fields:earth}{deklinasinya} $\ang{7}$ ke barat, sebuah kapal yang
dikemudikan ke ``utara kompas'' sejauh \qty{5.0}{km} hanyut
$5000 \times \sin\ang{7} \approx
\qty{6.1e2}{m}$ ke barat dari
jalurnya: untuk berlayar ke utara sejati, kemudikanlah $\ang{7}$ ke
\emph{timur} dari jarumnya.
\end{example}

\section{Medan magnet sebuah arus}

\begin{remark}[Oersted, 1820]\label{rem:g11:magnetic-fields:oersted}
Pada sebuah kuliah tahun 1820, Hans Christian Oersted melihat jarum
kompas berayun ketika ia menutup sebuah rangkaian di dekatnya:
kelistrikan dan kemagnetan ternyata satu pokok bahasan --- \emph{arus
menciptakan \omterm{def:g11:magnetic-fields:bfield}{medan magnet}}.
\end{remark}

\begin{proposition}[Medan sebuah kawat lurus]\label{prop:g11:magnetic-fields:wire}
Sebuah kawat lurus yang panjang dan dialiri arus $I$ menciptakan \omterm{def:g11:electric-gravitational-fields:field}{medan}
yang garisnya berupa lingkaran berpusat pada kawat itu, pada bidang
yang tegak lurus kawatnya. Pada jarak $d$ dari kawatnya,
\[
B = \frac{\mu_0\, I}{2\pi d},
\qquad \mu_0 = 4\pi\times 10^{-7}\ \unit{T.m/A} \approx \qty{1.26e-6}{T.m/A}.
\]
\end{proposition}
\admitted

\begin{omfigure}
\begin{tikzpicture}[scale=0.85]
  \node[font=\large] at (0,0) {$\odot$};
  \foreach \r in {0.55,1.05,1.55}
    \draw[omThm,-Stealth] (\r,0) arc (0:350:\r);
  \node[omThm, font=\small] at (45:1.9) {$\vect B$};
  \node[font=\small] at (0,-2.1) {arus menuju kita};
\end{tikzpicture}\qquad\qquad
\begin{tikzpicture}[scale=0.85]
  \node[font=\large] at (0,0) {$\otimes$};
  \foreach \r in {0.55,1.05,1.55}
    \draw[omThm,-Stealth] (\r,0) arc (360:10:\r);
  \node[omThm, font=\small] at (45:1.9) {$\vect B$};
  \node[font=\small] at (0,-2.1) {arus menjauhi kita};
\end{tikzpicture}

{\small \omterm{def:g11:electric-gravitational-fields:lines}{Garis medan} yang melingkar di sekeliling kawat lurus, dilihat
dari depan ($\odot$: arusnya keluar dari halaman; $\otimes$: masuk ke
halaman). Genggamlah kawatnya dengan tangan kanan, ibu jari searah
arusnya: jari-jarinya melingkupi ke arah berputarnya $\vect B$.}
\end{omfigure}

\begin{example}[Kawat adalah magnet yang lemah]\label{ex:g11:magnetic-fields:wire}
Pada $d = \qty{1.0}{cm}$ dari kawat yang dialiri $I = \qty{15}{A}$:
$B = \num{2e-7} \times 15 / 0.010 = \qty{0.30}{mT}$ --- enam kali \omterm{def:g11:electric-gravitational-fields:field}{medan}
Bumi, dari arus yang sudah akan memutuskan pemutus rumah tangga. \omterm{def:g11:electric-gravitational-fields:field}{Medan}
yang kuat memerlukan bentuk susunan yang lebih baik.
\end{example}

\begin{definition}[Kumparan dan solenoida]\label{def:g11:magnetic-fields:solenoid}
Menggulung kawatnya memusatkan \omterm{def:g11:electric-gravitational-fields:field}{medannya}. \emph{Kumparan
datar}\index{kumparan} ($N$ lilitan dalam satu cakram) berlaku seperti
\omterm{def:g11:magnetic-fields:magnet}{magnet} yang \omterm{def:g11:lenses-and-eye:lens}{tipis}, satu mukanya utara dan mukanya yang lain selatan;
\emph{solenoida}\index{solenoida} --- gulungan silindris yang panjang,
dengan $n$ lilitan tiap \omterm{def:g10:orders-of-magnitude:unit}{meter} --- berlaku seperti \omterm{def:g11:magnetic-fields:magnet}{magnet} batang.
\end{definition}

\begin{proposition}[Medan di dalam solenoida]\label{prop:g11:magnetic-fields:solenoidfield}
Di dalam \omterm{def:g11:magnetic-fields:solenoid}{solenoida} yang panjang, jauh dari ujungnya, \omterm{def:g11:electric-gravitational-fields:field}{medannya}
\emph{seragam}: sejajar sumbunya dan sama pada setiap titik di
dalamnya, dengan besar $B = \mu_0\, n\, I$, yang tidak bergantung pada
lebar rongganya. Di luarnya, \omterm{def:g11:electric-gravitational-fields:field}{medannya} lemah.
\end{proposition}
\admitted

\begin{remark}
Kedua rumusnya diturunkan secara jujur pada jilid Tahun ke-1, dari
hukum yang mengaitkan sirkulasi sebuah \omterm{def:g11:electric-gravitational-fields:field}{medan} dengan arus yang
dilingkupinya. Perhatikanlah apa yang berperan: $n$, yaitu lilitan
\emph{tiap \omterm{def:g10:orders-of-magnitude:unit}{meter}} --- gulunglah lebih rapat, bukan lebih panjang.
\end{remark}

\begin{method}[Tangan kanan, dua kali]\label{met:g11:magnetic-fields:righthand}
\begin{itemize}
  \item \emph{Kawat}: ibu jari = arusnya; jari yang melingkup = arah
        berputarnya \omterm{def:g11:electric-gravitational-fields:lines}{garis medannya}.
  \item \emph{\omterm{def:g11:magnetic-fields:solenoid}{Kumparan}, \omterm{def:g11:magnetic-fields:solenoid}{solenoida}}: jari yang melingkup = arus pada
        lilitannya; ibu jari = $\vect B$ di dalamnya, yang menunjuk ke
        muka utaranya.
\end{itemize}
\end{method}

\begin{omfigure}
\begin{tikzpicture}[scale=0.9]
  \foreach \x in {0,0.65,...,3.9} {
    \node[font=\small] at (\x,0.85) {$\odot$};
    \node[font=\small] at (\x,-0.85) {$\otimes$};}
  \draw[-Stealth, omThm, thick] (-0.4,0.4) -- (4.3,0.4);
  \draw[-Stealth, omThm, thick] (-0.4,0) -- (4.3,0)
    node[right, font=\small] {$\vect B$};
  \draw[-Stealth, omThm, thick] (-0.4,-0.4) -- (4.3,-0.4);
  \draw[omExo, -Stealth, dashed]
    (4.5,0.7) .. controls (5.9,2.3) and (-2.0,2.3) .. (-0.6,0.7);
\end{tikzpicture}

{\small Sebuah \omterm{def:g11:magnetic-fields:solenoid}{solenoida} yang dipotong memanjang ($\odot$/$\otimes$:
lilitan yang menyeberangi halaman): di dalamnya, $B = \mu_0 n I$ yang
seragam sepanjang sumbunya; di luarnya, \omterm{def:g11:electric-gravitational-fields:field}{medan} balik yang lemah dan
terbentang (garis putus-putus).}
\end{omfigure}

\begin{example}[Solenoida dalam angka]\label{ex:g11:magnetic-fields:solenoid}
Pada $10$ lilitan tiap sentimeter ($n = \qty{1000}{m^{-1}}$) dan
$I = \qty{5.0}{A}$: $B = \num{1.26e-6} \times 1000 \times 5.0 \approx
\qty{6.3}{mT}$ --- seratus kali \omterm{def:g11:electric-gravitational-fields:field}{medan} Bumi, yang dapat diatur dengan
sebuah kenop.
\end{example}

\section{Elektromagnet}

\begin{definition}[Elektromagnet]\label{def:g11:magnetic-fields:electromagnet}
Sebuah \emph{elektromagnet}\index{elektromagnet} adalah \omterm{def:g11:magnetic-fields:solenoid}{solenoida} yang
digulung pada inti besi lunak: besinya termagnetkan di dalam \omterm{def:g11:electric-gravitational-fields:field}{medan}
\omterm{def:g11:magnetic-fields:solenoid}{kumparannya} lalu melipatgandakannya, biasanya sebesar $100$ sampai
$1000$ kali. Tidak seperti \omterm{def:g11:magnetic-fields:magnet}{magnet} tetap, elektromagnet mati bersama
arusnya.
\end{definition}

\begin{remark}[Tiga mesin]\label{rem:g11:magnetic-fields:machines}
Derek di tempat rongsokan mengangkat sebuah mobil dengan \omterm{def:g11:magnetic-fields:electromagnet}{elektromagnet}
dan --- inilah intinya --- \emph{menjatuhkannya} dengan membuka sebuah
sakelar. \omterm{def:g11:magnetic-fields:electromagnet}{Elektromagnet} kecil sebuah relai menarik sebuah kontak sampai
tertutup: arus yang lemah menyalakan arus yang kuat
(\cref{ch:g11:circuits-and-power}). Pemindai MRI menahan
\qty{1.5}{T} yang seragam di sekeliling pasiennya dengan \omterm{def:g11:magnetic-fields:solenoid}{solenoida}
adihantar.
\end{remark}

\omterm{def:g10:orders-of-magnitude:oom}{Orde besarnya}, dari galaksi sampai bintang yang mati:
\begin{center}
\small
\begin{tabular}{lr}
ruang antarbintang & $\sim \qty{e-10}{T}$ \\
\omterm{def:g11:electric-gravitational-fields:field}{medan} permukaan Bumi & \qty{5e-5}{T} \\
\omterm{def:g11:magnetic-fields:magnet}{magnet} lemari es & $\sim \qty{5e-3}{T}$ \\
\omterm{def:g11:magnetic-fields:magnet}{magnet} neodimium, pada kutubnya & $\sim \qty{0.5}{T}$ \\
\omterm{def:g11:magnetic-fields:solenoid}{solenoida} MRI & \qty{1.5}{T} \\
permukaan bintang neutron & $\sim \qty{e8}{T}$ \\
\end{tabular}
\end{center}

\section{Gaya Laplace}

\begin{proposition}[Gaya Laplace]\label{prop:g11:magnetic-fields:laplace}
Sebuah kawat lurus sepanjang $L$, yang dialiri arus $I$ di dalam \omterm{def:g11:electric-gravitational-fields:capacitor}{medan
seragam} $\vect B$ yang \emph{tegak lurus} kawatnya, merasakan
\emph{gaya Laplace}\index{gaya Laplace} sebesar $F = I\,L\,B$, tegak
lurus terhadap kawatnya sekaligus terhadap $\vect B$: bukalah telapak
tangan kanan, ibu jari searah arusnya, jari yang lurus searah $\vect B$
--- telapaknya mendorong searah $\vect F$. Kawat yang \emph{sejajar}
\omterm{def:g11:electric-gravitational-fields:field}{medannya} tidak merasakan \omterm{def:g10:inertia:force}{gaya}.
\end{proposition}
\admitted

\begin{remark}[Rel yang meloncat]\label{rem:g11:magnetic-fields:rail}
Jilid Tahun ke-1 menurunkannya dari \omterm{def:g10:inertia:force}{gaya} \omterm{def:g11:magnetic-fields:magnet}{magnet} pada tiap muatan yang
bergerak. Di laboratorium: sebatang tembaga terletak melintang pada dua
rel mendatar di antara kutub sebuah \omterm{def:g11:magnetic-fields:magnet}{magnet} (\omterm{def:g11:electric-gravitational-fields:field}{medannya} masuk ke halaman
pada gambar); tutuplah sakelarnya, dan batang itu meluncur sepanjang
relnya. Baliklah arusnya, atau baliklah \omterm{def:g11:magnetic-fields:magnet}{magnetnya}: batang itu meluncur
ke arah yang lain.
\end{remark}

\begin{omfigure}
\begin{tikzpicture}[scale=0.95]
  \foreach \x in {0.6,1.5,2.9,3.8}
    \foreach \y in {0.35,1.05}
      \node[omThm, font=\small] at (\x,\y) {$\otimes$};
  \node[omThm, font=\small] at (4.35,1.05) {$\vect B$};
  \draw[very thick] (0,0) -- (4.6,0);
  \draw[very thick] (0,1.4) -- (4.6,1.4);
  \draw[omDef, line width=2.5pt] (2.1,-0.12) -- (2.1,1.52);
  \draw[-Stealth, omDef, thick] (0.1,1.6) -- (1.6,1.6)
    node[midway, above, font=\small] {$I$};
  \draw[-Stealth, omDef, thick] (2.32,1.15) -- (2.32,0.45);
  \draw[-Stealth, omDef, thick] (1.6,-0.22) -- (0.1,-0.22);
  \draw[-Stealth, omMeth, very thick] (2.1,0.7) -- (3.4,0.7)
    node[pos=0.85, above, font=\small] {$\vect F$};
\end{tikzpicture}

{\small Percobaan rel: arusnya turun menyusuri batangnya, \omterm{def:g11:electric-gravitational-fields:field}{medannya}
masuk ke halaman ($\otimes$) --- \omterm{prop:g11:magnetic-fields:laplace}{gaya Laplace} $F = ILB$ menggerakkan
batang itu sepanjang relnya.}
\end{omfigure}

\begin{example}[Laplace dalam angka]\label{ex:g11:magnetic-fields:laplace}
Sebatang sepanjang $L = \qty{5.0}{cm}$ yang dialiri $I = \qty{8.0}{A}$
melintasi $B = \qty{0.50}{T}$:
$F = 8.0 \times 0.050 \times 0.50 = \qty{0.20}{N}$ --- setara \omterm{def:g10:universal-gravitation:weight}{berat}
\qty{20}{g}, pada kawat yang beratnya beberapa gram.
\end{example}

\begin{remark}[Motor arus searah]\label{rem:g11:magnetic-fields:motor}
Tekuklah kawatnya menjadi sebuah untai di dalam \omterm{def:g11:electric-gravitational-fields:field}{medannya}: kedua sisinya
yang tegak lurus $\vect B$ membawa arus yang berlawanan, sehingga \omterm{prop:g11:magnetic-fields:laplace}{gaya
Laplace-nya} berlawanan --- sepasang \omterm{def:g10:inertia:force}{gaya} yang memutar untainya.
Setengah putaran kemudian pasangan itu akan memutarnya kembali,
sehingga sebuah sakelar berputar (\emph{komutator}\index{komutator})
membalikkan arusnya tiap setengah putaran. Setiap kipas dan setiap bor
adalah untai ini, yang dilipatgandakan.
\end{remark}

\section{Latihan}

\begin{exercise}[$\star$]\label{exo:g11:magnetic-fields:1}
Sebuah \omterm{def:g11:magnetic-fields:magnet}{magnet} batang digergaji menjadi dua di antara kutubnya. Kutub
apa yang dibawa tiap potongnya? Jika didekatkan kembali, apakah muka
yang baru terpotong itu tarik-menarik atau tolak-menolak? Adakah cara
pemotongan yang dapat mengasingkan kutub utaranya?
\end{exercise}

\begin{exercise}[$\star$]\label{exo:g11:magnetic-fields:2}
Benar atau salah, masing-masing dengan satu alasan: (a) dua \omterm{def:g11:magnetic-fields:lines}{garis medan
magnet} berpotongan di tempat \omterm{def:g11:electric-gravitational-fields:field}{medannya} kuat; (b) di luar sebuah \omterm{def:g11:magnetic-fields:magnet}{magnet},
\omterm{def:g11:electric-gravitational-fields:lines}{garis medannya} berjalan dari kutub utara ke kutub selatan; (c) setiap
\omterm{def:g11:magnetic-fields:lines}{garis medan magnet} berupa untai tertutup; (d) sebuah jarum berhenti
tegak lurus terhadap garis yang melewatinya.
\end{exercise}

\begin{exercise}[$\star$]\label{exo:g11:magnetic-fields:3}
Hitunglah \omterm{def:g11:electric-gravitational-fields:field}{medan} \qty{2.0}{cm} dari kawat lurus yang dialiri
\qty{10}{A}. Berapa kali \omterm{def:g11:electric-gravitational-fields:field}{medan} Bumi yang \qty{50}{\micro T} itu?
\end{exercise}

\begin{exercise}[$\star$]\label{exo:g11:magnetic-fields:4}
Sebuah \omterm{def:g11:magnetic-fields:solenoid}{solenoida} berlilitan $800$ yang digulung sepanjang \qty{40}{cm}
dialiri $I = \qty{1.5}{A}$. Hitunglah $n$, lalu $B$ di dalamnya;
menjadi berapa $B$ jika arusnya digandakan?
\end{exercise}

\begin{exercise}[$\star$]\label{exo:g11:magnetic-fields:5}
Sepotong kawat sepanjang \qty{10}{cm} dialiri \qty{5.0}{A} tegak lurus
\omterm{def:g11:electric-gravitational-fields:field}{medan} \qty{0.20}{T}. Hitunglah \omterm{prop:g11:magnetic-fields:laplace}{gaya Laplace-nya}; \omterm{def:g10:universal-gravitation:weight}{berat} massa berapa
yang setara dengan \omterm{def:g10:inertia:force}{gaya} itu?
\end{exercise}

\begin{exercise}[$\star\star$]\label{exo:g11:magnetic-fields:6}
Di tempat seorang pendaki berdiri, \omterm{def:g11:magnetic-fields:earth}{deklinasinya} $\ang{8}$ ke barat. Ia
berjalan \qty{2.0}{km} ``ke utara menurut kompas'': sejauh apa, dan ke
sisi mana dari utara sejati, ia hanyut? Arah mana yang seharusnya
membawanya ke utara sejati?
\end{exercise}

\begin{exercise}[$\star\star$]\label{exo:g11:magnetic-fields:7}
Pada jarak berapa dari kawat lurus yang dialiri \qty{20}{A} \omterm{def:g11:electric-gravitational-fields:field}{medannya}
menyamai \omterm{def:g11:electric-gravitational-fields:field}{medan} Bumi yang \qty{50}{\micro T}? Apa yang dikatakan hal itu
tentang memercayai kompas di dekat kabel yang bertegangan?
\end{exercise}

\begin{exercise}[$\star\star$]\label{exo:g11:magnetic-fields:8}
Rancanglah \omterm{def:g11:magnetic-fields:solenoid}{solenoida} yang menghasilkan $B = \qty{10}{mT}$ dengan
$I = \qty{4.0}{A}$: hitunglah $n$ yang diperlukan, lalu banyaknya
lilitan untuk \omterm{def:g11:magnetic-fields:solenoid}{kumparan} sepanjang \qty{25}{cm}.
\end{exercise}

\begin{exercise}[$\star\star$]\label{exo:g11:magnetic-fields:9}
Pada percobaan rel, batangnya (bermassa \qty{20}{g}, panjang
$L = \qty{12}{cm}$ di antara relnya) dialiri \qty{10}{A} di dalam \omterm{def:g11:electric-gravitational-fields:field}{medan}
\qty{0.50}{T} yang tegak lurus. Hitunglah \omterm{prop:g11:magnetic-fields:laplace}{gaya Laplace-nya}, lalu
percepatan awal batangnya; bandingkan dengan $g$.
\end{exercise}

\begin{exercise}[$\star\star$]\label{exo:g11:magnetic-fields:10}
Sebuah kawat mendatar bermassa \qty{5.0}{g} dan panjang \qty{20}{cm}
berada di dalam \omterm{def:g11:electric-gravitational-fields:field}{medan} mendatar $B = \qty{0.10}{T}$ yang tegak lurus
terhadapnya. Arus berapa yang membuat \omterm{prop:g11:magnetic-fields:laplace}{gaya Laplace-nya} mengimbangi
\omterm{def:g10:universal-gravitation:weight}{berat} kawat itu, sehingga kawatnya melayang? Ke arah mana \omterm{def:g10:inertia:force}{gayanya} harus
menunjuk?
\end{exercise}

\begin{exercise}[$\star\star$]\label{exo:g11:magnetic-fields:11}
Sebutkan arah \omterm{prop:g11:magnetic-fields:laplace}{gaya Laplace-nya} (atau nyatakan bahwa \omterm{def:g10:inertia:force}{gayanya} lenyap):
(a) arusnya ke kanan halaman, $\vect B$ masuk ke halaman;
(b) arusnya ke arah atas halaman, $\vect B$ keluar dari halaman;
(c) arusnya sejajar $\vect B$. Pakailah telapak tangan kanan yang
terbuka.
\end{exercise}

\begin{exercise}[$\star\star\star$]\label{exo:g11:magnetic-fields:12}
Sebuah kawat mendatar terbentang utara--selatan secara geografis,
\qty{3.0}{cm} di atas sebuah kompas; komponen mendatar \omterm{def:g11:electric-gravitational-fields:field}{medan} Bumi di
sana \qty{20}{\micro T}. Kawat itu dialiri $I = \qty{5.0}{A}$:
hitunglah \omterm{def:g11:electric-gravitational-fields:field}{medannya} di kompas itu, sebutkan arahnya, lalu simpangan
jarumnya (\omterm{def:g11:electric-gravitational-fields:field}{medan} berjumlah sebagai vektor).
\end{exercise}

\begin{exercise}[$\star\star\star$]\label{exo:g11:magnetic-fields:13}
Sebuah \omterm{def:g11:magnetic-fields:electromagnet}{elektromagnet} tempat rongsokan berupa \omterm{def:g11:magnetic-fields:solenoid}{kumparan} berlilitan $400$
sepanjang \qty{20}{cm}, berhambatan \qty{2.0}{\ohm}, pada pasokan
\qty{12}{V}; inti besinya melipatgandakan \omterm{def:g11:electric-gravitational-fields:field}{medan} telanjangnya sebesar
$100$ kali. Hitunglah arusnya, \omterm{def:g11:electric-gravitational-fields:field}{medan} \omterm{def:g11:magnetic-fields:solenoid}{kumparan} telanjangnya, \omterm{def:g11:electric-gravitational-fields:field}{medannya}
beserta intinya, dan \omterm{def:g10:energy-conservation:power}{daya} yang dibuang \omterm{def:g11:magnetic-fields:solenoid}{kumparannya}
(\cref{ch:g11:circuits-and-power}). Mengapa rongsokannya jatuh seketika
ketika sakelarnya dibuka?
\end{exercise}

\begin{exercise}[$\star\star\star$]\label{exo:g11:magnetic-fields:14}
Untai persegi panjang sebuah motor memiliki $N = 50$ lilitan yang
dialiri $I = \qty{2.0}{A}$; kedua sisinya yang panjangnya \qty{5.0}{cm}
tegak lurus \omterm{def:g11:electric-gravitational-fields:field}{medan} $B = \qty{0.30}{T}$. Hitunglah \omterm{def:g10:inertia:force}{gaya} pada
masing-masing sisi itu; mengapa kedua \omterm{def:g10:inertia:force}{gayanya} berlawanan, dan apa yang
dilakukan keduanya pada untainya? Mengapa kedua sisinya yang lain
kadang-kadang tidak merasakan \omterm{def:g10:inertia:force}{gaya}? Mengapa arusnya harus dibalik tiap
setengah putaran?
\end{exercise}

\begin{exercise}[$\star\star\star$]\label{exo:g11:magnetic-fields:15}
Sedekat apa dengan kabel yang dialiri \qty{100}{A} kamu harus berada
agar \omterm{def:g11:electric-gravitational-fields:field}{medannya} mencapai \qty{1}{T}? Bandingkan dengan jari-jari kabelnya
yang beberapa milimeter, lalu simpulkan mengapa \omterm{def:g11:electric-gravitational-fields:field}{medan} sekuat \omterm{def:g11:magnetic-fields:bfield}{tesla}
dibangun dari \omterm{def:g11:magnetic-fields:solenoid}{solenoida}, bukan dari kawat tunggal. Lalu: \omterm{def:g11:electric-gravitational-fields:field}{medan}
permukaan bintang neutron sebesar \qty{e8}{T} berapa pangkat sepuluh di
atas \omterm{def:g11:electric-gravitational-fields:field}{medan} Bumi?
\end{exercise}

\section{Soal: Kompas, derek, dan MRI}

\begin{problem}[{Soal akhir pekan --- tiga magnet yang bekerja: sebuah
kompas menyeberangi samudra, sebuah elektromagnet mengangkat mobil, dan
sebuah solenoida MRI --- satu medan, membentang sepuluh pangkat
sepuluh, mengerjakan tiga pekerjaan}]
\label{pb:g11:magnetic-fields:1}
Vektor $\vect B$ yang sama mengemudikan perahu layar dengan
\qty{50}{\micro T}, mengangkat besi rongsokan dengan \qty{1}{T}, dan
mencitrakan sebuah lutut dengan \qty{1.5}{T} --- tiga mesin untuk
ketiga rumus bab ini.

\medskip
\noindent\textbf{Bagian I --- Kompas.}
Pada kedudukan kapalnya, \omterm{def:g11:electric-gravitational-fields:field}{medan} Bumi memiliki komponen mendatar
\qty{20}{\micro T} dan komponen tegak \qty{44}{\micro T};
\omterm{def:g11:magnetic-fields:earth}{deklinasinya} $\ang{6}$ ke barat.
\begin{enumerate}
  \item Mengapa jarum kompas menunjuk ke utara (\omterm{def:g11:magnetic-fields:magnet}{magnet}) sama sekali?
  \item Hitunglah besar \omterm{def:g11:electric-gravitational-fields:field}{medan} seluruhnya dan sudutnya terhadap arah
        mendatar.
  \item Dengan dikemudikan ke ``utara kompas'' sejauh \qty{30}{km},
        berapa jauh ke barat dari jalur utara sejatinya kapal itu
        berakhir?
  \item Muatan bajanya menyimpangkan jarumnya $\ang{4}$ lagi ke barat:
        pertanyaan yang sama.
  \item Di dekat kutub geomagnetnya sebuah kompas tidak berguna:
        komponen yang mana yang harus disalahkan, dan mengapa?
\end{enumerate}

\medskip
\noindent\textbf{Bagian II --- Dereknya.}
\omterm{def:g11:magnetic-fields:electromagnet}{Elektromagnet} tempat rongsokan itu berupa \omterm{def:g11:magnetic-fields:solenoid}{solenoida} berlilitan $600$
yang digulung sepanjang \qty{30}{cm}, berhambatan \qty{1.5}{\ohm}, yang
dipasok \qty{12}{V}; inti besinya melipatgandakan \omterm{def:g11:electric-gravitational-fields:field}{medan} telanjangnya
sebesar $50$ kali.
\begin{enumerate}[resume]
  \item Hitunglah $n$, yaitu kerapatan lilitannya.
  \item Hitunglah arus pada \omterm{def:g11:magnetic-fields:solenoid}{kumparannya}
        (\cref{ch:g11:circuits-and-power}).
  \item Hitunglah \omterm{def:g11:electric-gravitational-fields:field}{medan} \omterm{def:g11:magnetic-fields:solenoid}{kumparan} telanjangnya $\mu_0 n I$.
  \item Hitunglah \omterm{def:g11:electric-gravitational-fields:field}{medannya} beserta intinya; bandingkan dengan
        \qty{0.5}{T} milik neodimium.
  \item Hitunglah \omterm{def:g10:energy-conservation:power}{daya} yang dibuang \omterm{def:g11:magnetic-fields:solenoid}{kumparannya}, lalu \omterm{def:g10:energy-conservation:energy}{energinya} untuk
        satu giliran mengangkat selama \qty{10}{\minute}.
  \item \omterm{def:g11:magnetic-fields:magnet}{Magnet} tetap sekuat itu memang ada. Lalu mengapa tetap memakai
        \omterm{def:g11:magnetic-fields:electromagnet}{elektromagnet}?
\end{enumerate}

\medskip
\noindent\textbf{Bagian III --- Motor derek gulungnya.}
Derek gulung itu digerakkan motor arus searah: sebuah untai persegi
panjang berlilitan $N = 100$, dengan sisi sepanjang \qty{4.0}{cm} yang
tegak lurus $B = \qty{0.25}{T}$, dialiri $I = \qty{3.0}{A}$.
\begin{enumerate}[resume]
  \item Hitunglah \omterm{prop:g11:magnetic-fields:laplace}{gaya Laplace} pada tiap berkas sisi yang tegak lurus
        $\vect B$.
  \item Kedua \omterm{def:g10:inertia:force}{gayanya} sama besar dan berlawanan, namun akibatnya tidak
        saling menghapus. Apa yang dihasilkan keduanya, dan mengapa
        keduanya harus \emph{berlawanan}?
  \item \omterm{def:g10:inertia:force}{Gaya} apa yang bekerja pada kedua sisi yang lain, ketika sisi
        itu sejajar $\vect B$?
  \item Setelah setengah putaran, \omterm{def:g10:inertia:force}{gaya} yang sama akan membatalkan
        putarannya. Apa yang dikerjakan \omterm{rem:g11:magnetic-fields:motor}{komutatornya}, dan kapan?
  \item Sebutkan tiga perubahan rancangan yang berbeda, yang
        masing-masing menggandakan pasangan \omterm{def:g10:inertia:force}{gayanya}.
\end{enumerate}

\medskip
\noindent\textbf{Bagian IV --- MRI.}
\omterm{def:g11:magnetic-fields:solenoid}{Solenoida} pemindai itu menahan \qty{1.5}{T} yang seragam dengan arus
$I = \qty{500}{A}$.
\begin{enumerate}[resume]
  \item Berapa kali \omterm{def:g11:electric-gravitational-fields:field}{medan} Bumi nilai \qty{1.5}{T} itu?
  \item Hitunglah kerapatan lilitan $n$ yang diperlukan \emph{tanpa}
        inti besi apa pun.
  \item Jika gulungannya berhambatan \qty{0.20}{\ohm}, \omterm{def:g10:energy-conservation:power}{daya} berapa yang
        akan dibuangnya? Sifat apa pada gulungan yang sebenarnya yang
        menghindarkan hal itu, dengan harga kedinginan yang luar biasa?
  \item Puncak ceritanya: tempatkan \omterm{def:g11:electric-gravitational-fields:field}{medan} kapalnya, dereknya dan
        MRI-nya pada tangga pangkat sepuluh dari ruang antarbintang
        (\qty{e-10}{T}) sampai bintang neutron (\qty{e8}{T}); apa yang
        berubah di antara ketiga mesin itu --- fisikanya, ataukah
        \omterm{def:g11:circuits-and-power:current}{amperenya}?
\end{enumerate}
\end{problem}
